\subsection{Surveying Ethical Theories for AI Alignment}
\label{sec:2.1.1}
Six frameworks do work in this book. Each is given here with the case it handles well and the case that defeats it, because the defeats are what force the structure chapter~\ref{sec:3} ends up proposing.

\begin{enumerate}
  \item Consequentialist theories, such as utilitarianism, judge actions by their ramifications rather than their nature, aiming to maximize well-being and minimize harm. Applied to AI, the objective becomes optimizing human welfare — consider an AI system designed to distribute scarce medical resources in a crisis, prioritizing treatment by predicted outcomes for patient health and life expectancy. The complexity comes from forecasting those long-term consequences accurately, which can put individual rights and collective welfare into direct conflict.
  \item Deontological theories weigh adherence to moral rules and duties over consequences. An AI system built on deontological principles would hold to specific obligations — respecting privacy, avoiding harm, truthful communication — regardless of outcome: a conversational AI system programmed never to deceive users even when deception would yield a short-term benefit. The challenge is resolving conflicts between duties, or handling cases the rules simply do not cover.
  \item Virtue ethics focuses on developing moral character — traits like honesty, courage, and empathy — rather than following rules or calculating outcomes. An AI system built on virtue ethics would aim to embody and encourage those traits in its interactions: a personal-assistant AI system that models honesty consistently in its own responses. The difficulty is identifying which virtues fit a given context and translating an intangible trait like courage into a concrete action.
  \item Care ethics centers relationships and responsiveness to another's needs. An AI system built on care ethics would prioritize nurturing those bonds and responding well to what a person actually needs — an elder-care robot attending not just to physical needs but to company and emotional support. Section~\ref{sec:2.2} takes up why the acting half of that, compassion, and the feeling half, empathy, are not the same capacity, and why an AI system needs the first far more than the second. The challenge here is tending to several stakeholders' different, sometimes conflicting needs at once.
  \item Contractualism holds that moral principles should come from agreements among rational individuals. An AI system guided by contractualism would try to uphold whatever principles were mutually agreed — a system managing a shared resource like an energy grid, allocating it by pre-agreed priorities. The hurdle is accommodating diverse perspectives and reaching consensus on shared values in the first place.
  \item The capabilities approach asks not what an AI system does but what it leaves a person able to be and do. Applied to AI, the test of a system is whether it widens the range of lives actually available to the people it touches: an adaptive tutor that meets a student where they are expands a capability; the same system gating what a student may study expands nothing, however well it teaches inside the gate. The difficulty is that capability is harder to measure than either an outcome or a rule, and a system optimizing something easier to measure will drift toward that instead. Section~\ref{sec:10.1} takes this framework up in full.
\end{enumerate}
Two of the six are load-bearing later and the rest are equipment. The capabilities approach is what chapter~\ref{sec:10} uses to ask what a life touched by an AI system is actually free to become, which is a different question from whether the system behaved well on any particular occasion. Deontology is the one that returns: its unconditional constraints are what chapter~\ref{sec:3} needs, and the objection that sinks it as a complete theory — duties conflict, rules run out — has no force against a constraint that decides a small number of things and leaves the rest to the other five.

That is the pluralism this section is actually recommending, and it is not the usual one. The usual version blends frameworks case by case and hopes the blend is wiser than any of its parts, which leaves a system that can be argued out of anything, because for any act there is a framework in the blend that permits it. What section~\ref{sec:2.1.3} takes up is the mechanics of conflict among them; what chapter~\ref{sec:3} adds is that the conflict has to bottom out somewhere the arguing does not reach.
